We asked whether a model's explanations drift before its accuracy does.
Under our single-window rule they did, by up to 56 days. The signal
is fragile, however. When two windows in a row are required, the blind
model's $\tau_w$ never warns while its error does, and $\tau_w$ measured
against the same fortnight a year earlier never crosses for either model,
suggesting much of the early movement was seasonal. Rank-based explanation drift is therefore not yet a reliable early warning
here. The explanations changed more in scale than in order. Attribution to
yesterday's price tripled while the same four features stayed on top, growth
that a rank-based measure cannot register.

Our study has clear limits. It covers one event, one market and one model type (gradient-boosted
trees), so the results may not generalise. The models are frozen, so we
cannot say how monitoring interacts with retraining. The thresholds rest on only eight baseline windows. We used only TreeSHAP
with a mean-$|$SHAP$|$ ranking. Other explainers may differ.
Future work should test magnitude-aware drift measures, use seasonally
matched references by default, and study monitoring alongside retraining.